% Generated draft: requires mathematical and source review.
\documentclass[11pt]{article}
\usepackage[margin=25mm]{geometry}
\usepackage{fontspec}
\setmainfont{Latin Modern Roman}
\usepackage{amsmath,amssymb,mathtools,unicode-math}
\setmathfont{Latin Modern Math}
\usepackage{enumitem,needspace,xurl,hyperref}
\setlength{\parindent}{0pt}
\setlength{\parskip}{5pt}
\setlength{\emergencystretch}{3em}
\newcommand{\N}{\mathbb{N}}
\newcommand{\Z}{\mathbb{Z}}
\newcommand{\Q}{\mathbb{Q}}
\newcommand{\R}{\mathbb{R}}
\newcommand{\C}{\mathbb{C}}
\newcommand{\F}{\mathbb{F}}
\newcommand{\A}{\mathbb{A}}
\newcommand{\PP}{\mathbb{P}}
\newcommand{\E}{\mathbb{E}}
\newcommand{\eps}{\varepsilon}
\newcommand{\dd}{\,\mathrm{d}}
\begin{document}
\section*{UnsolvedMath ID: 30004512}
\textbf{Source problem code:} catalogue entry
\subsection{Definitions and mathematical statement}
The supplied record poses a research programme rather than a single quantified conjecture. It does not specify an ambient space, a class of admissible nonconvex families, or particular numerical bounds. The following definitions provide conventional language for its requested conclusions; they do not supply the missing hypotheses.

Let \(X\) be a set and \(\mathcal F\subseteq 2^X\) a family of subsets. The intersection-generated closure associated with \(\mathcal F\) is
\[
\operatorname{cl}_{\mathcal F}(A)=\bigcap\{F\in\mathcal F:A\subseteq F\},\qquad A\subseteq X,
\]
where an intersection over no sets is \(X\). More generally, a closure operator \(\operatorname{cl}:2^X\to2^X\) satisfies, for all \(A,B\subseteq X\),
\[
A\subseteq\operatorname{cl}(A),\qquad
A\subseteq B\Longrightarrow\operatorname{cl}(A)\subseteq\operatorname{cl}(B),\qquad
\operatorname{cl}(\operatorname{cl}(A))=\operatorname{cl}(A).
\]
The source does not specify whether the intended family is \(\mathcal F\), its intersection closure, or the closed sets of some separately given operator. Nor does it specify the ambient convexity relative to which sets are nonconvex.

A finite Helly bound \(h\geq1\) for \(\mathcal F\) means that every finite subfamily \(\mathcal A\subseteq\mathcal F\) satisfies
\[
\left[\bigcap\mathcal B\ne\varnothing\text{ for every }\mathcal B\subseteq\mathcal A\text{ with }1\leq|\mathcal B|\leq h\right]
\Longrightarrow\bigcap\mathcal A\ne\varnothing.
\]
For a closure operator, a Radon conclusion for a finite set \(P\subseteq X\) is the existence of a partition \(P=P_1\sqcup P_2\) into nonempty parts with
\[
\operatorname{cl}(P_1)\cap\operatorname{cl}(P_2)\ne\varnothing.
\]
For an integer \(r\geq2\), a Tverberg conclusion is a partition \(P=P_1\sqcup\cdots\sqcup P_r\) into nonempty parts such that
\[
\bigcap_{i=1}^{r}\operatorname{cl}(P_i)\ne\varnothing.
\]
Corresponding quantitative questions seek cardinality thresholds guaranteeing these conclusions for every admissible finite point set. Whether points may repeat, and whether conclusions must instead hold for sequences or multisets, is unspecified.

A conventional fractional Helly conclusion has an integer parameter \(k\geq1\): for each \(\alpha\in(0,1]\), there is \(\beta\in(0,1]\) such that, for every finite \(\mathcal A\subseteq\mathcal F\) with \(n=|\mathcal A|\geq k\),
\[
\left|\left\{\mathcal B\subseteq\mathcal A:|\mathcal B|=k,\ \bigcap\mathcal B\ne\varnothing\right\}\right|
\geq\alpha\binom nk
\Longrightarrow
\exists x\in X:\ \left|\{A\in\mathcal A:x\in A\}\right|\geq\beta n.
\]
A conventional weak \(\varepsilon\)-net conclusion, for \(0<\varepsilon\leq1\) and a finite nonempty \(P\subseteq X\), seeks a finite \(N\subseteq X\) meeting every \(F\in\mathcal F\) with \(|F\cap P|\geq\varepsilon|P|\), together with a bound on \(|N|\) independent of \(|P|\). Here the net points need not belong to \(P\). The source does not select this finite uniform formulation over possible measure-based variants or specify the threshold convention.

For integers \(p\geq q\geq2\), a finite family \(\mathcal A\subseteq\mathcal F\) with \(|\mathcal A|\geq p\) has the \((p,q)\)-property if every \(p\)-member subfamily contains \(q\) members with nonempty common intersection. A corresponding piercing conclusion seeks a finite \(T\subseteq X\) meeting every member of \(\mathcal A\), with a uniform bound on \(|T|\).

\textbf{Source question.} Determine hypotheses on nonconvex set systems or their associated closure operators under which Helly, Radon, Tverberg, fractional Helly, weak \(\varepsilon\)-net, and \((p,q)\)-type conclusions hold, and obtain quantitative bounds with explicit parameter dependence. The admissible systems, hypotheses, dependencies, and intended variants remain unspecified. Consequently the supplied question cannot be sharpened to one definite theorem or conjecture without additional information. The supplied background mentions Betti-number and topological-complexity hypotheses, but gives no precise assumptions or bounds to formalize.
\subsection{Short English statement}
This is a broad research question about extending intersection, partition, and piercing theorems to nonconvex families and closure systems. The definitions describe conventional possible targets, but the source leaves the ambient setting, hypotheses, variants, and quantitative bounds unspecified. Its supplied partially solved label is dated curation, not an independently verified mathematical conclusion.
\subsection{Sources}
\begin{itemize}
\item Local registry: \texttt{data/unsolvedmath/problems.json}, permanent UnsolvedMath ID \(30004512\), source code \texttt{catalogue entry}, entitled ``Helly--Radon--Tverberg Theory for Nonconvex Systems.'' The supplied collection is Oberwolfach Reports Open Problems; no individual proposer is named. The supplied proposed year is \(2020\).
\item Supplied report citation: \emph{Discrete Geometry (2021)}, Oberwolfach Reports, DOI 10.4171/owr/2020/30, pp.\ 1469--1528. \url{https://doi.org/10.4171/owr/2020/30}. The background attributes the question to mathematical workshop open-problem discussions documented by Oberwolfach Reports; no precise problem location within those pages is supplied.
\item Supplied literature source: Patakova, \emph{Bounding Radon Numbers via Betti Numbers}. \url{https://doi.org/10.4230/LIPIcs.SoCG.2020.61}.
\item Supplied literature source: Patakova, preprint. \url{https://arxiv.org/abs/1908.01677}.
\item Literature assessment dated exactly \texttt{2026-08-22}: the supplied status is ``partially solved.'' The supplied assessment states that Betti-number and topological-complexity hypotheses yield Radon and Helly-type bounds for many nonconvex families, while the unrestricted question remains open. This is supplied dated curation derived from the OWR material, not independently verified here; its time-sensitive status and cited claims require independent checking.
\item The registry labels the statement \texttt{corrected\_verified}, but its accompanying verification description limits the repairs to formula delimiters, notation, grammar, and typography while preserving mathematical content and scope. This label does not establish mathematical correctness or resolve the formulation gaps.
\end{itemize}
\end{document}
